%% =====================================================================
%%  T-02-01  Affect Control
%%  -------------------------------------------------------------------
%%  One idea per file. Core is mandatory. Slots in canonical order.
%%  Max 700 words. No layout commands. No filler. New source material
%%  for this entry goes in sources/<BOOK>/T-02-01.tex -- never by
%%  rewriting this file (docs/RULES.md R0.1).
%% =====================================================================
%% @kind: topic
%% @id: T-02-01
%% @title: Affect Control
%% @subtitle: Calm is a position, not a temperament
%% @chapter: c02
%% @order: 10
%% @status: stable
%% @sources: B3
%% @updated: 2026-09-19

\topic{T-02-01}{Affect Control}{Calm is a position, not a temperament}

\begin{core}
Whoever is more composed holds the exchange. Anger and excitement are strategically useless to the person feeling them and highly useful to the person provoking them, so the first piece of equipment is the ability to stay flat while the counterpart is not.
\end{core}
\begin{mechanism}
Strong affect narrows attention and shortens the horizon. An angry person optimises for the next thirty seconds and pays for it later; a calm person can hold the sequence in view. That asymmetry is exploitable in both directions, which is why the same source treats keeping yourself composed and making the opponent emotional as one technique rather than two. Composure also conceals information: a flat face gives away nothing about what you want, what you fear, or what you will accept.
\end{mechanism}
\begin{tells}
  \item Your voice or tempo changes when the subject turns to you.
  \item You rehearse a reply while the other person is still talking.
  \item You find yourself trying to win the exchange rather than to get the thing.
  \item Someone has located a subject that reliably rattles you, and keeps returning to it.
\end{tells}
\begin{moves}
  \item Identify in advance the two or three subjects that reliably upset you.
  \item Treat a rise in your own affect as a signal to stop advancing and start listening.
  \item Answer slowly and below the emotional register of the question.
  \item If the counterpart is escalating, let them; do not meet the register.
  \item Decide the outcome you want before the conversation, and judge each of your own moves against it rather than against the insult.
\end{moves}
\begin{counters}
  \item Recognise provocation as information: someone is trying to move you off your plan, which tells you the plan is working.
  \item Buy time. Delay is the cheapest affect control there is and almost nothing in this book survives it.
  \item Refuse to be the one who escalates in front of an audience; the audience will attribute the escalation to whoever is louder, not whoever started it.
  \item If you cannot stay flat, leave. A withdrawn exchange can be resumed; a lost one cannot.
\end{counters}
\begin{limits}
Flatness reads as coldness, and coldness costs trust in the relationships where trust is the currency. Applied everywhere it produces a person nobody confides in and therefore nobody tells anything to --- which destroys the information-gathering half of this book. It is a tool for contested exchanges, not a way of living, and it is also the single hardest item in this chapter to acquire, because it requires the self-knowledge listed in T-04-04 first.
\end{limits}

\sources{B3}
\seesources
\seealso{T-02-02, T-04-03, T-17-01, T-21-03}
